\bookchapter{Anhang A: “Babysitter”-Ablauf für ein neues Smartphone}

\booksection{Auspacken und Prüfen in drei Schritten}

\begin{itemize}
  \item \textbf{: Mit Apple-ID, Huawei-Konto, Xiaomi-Konto usw. anmelden.}Grundeinstellungen
  \item \textbf{: "Mein iPhone suchen" aktivieren und ein vernünftiges Sperrbildschirm-Passwort festlegen.}eSIM-Einrichtung
  \item \textbf{: Wenn dein Smartphone eSIM unterstützt (z. B. iPhone, einige Android-Flaggschiffe), kannst du in den Einstellungen ein eSIM-Tarifpaket hinzufügen, ohne eine physische SIM-Karte einzulegen. eSIM ist besonders beim Auslandsaufenthalt praktisch – du kannst direkt ein eSIM-Paket eines lokalen Anbieters kaufen.}Datenmigration
\end{itemize}

\booksection{Aktivierung und Ersteinrichtung – “Wichtige Erinnerungen”}

\begin{itemize}
  \item \textbf{3. App-Installation – "Sicherheit zuerst"}Zuerst der offizielle Store
  \item \textbf{: Apple nutzt den App Store, Android den vorinstallierten App-Store des Herstellers.}Berechtigungen nicht leichtfertig vergeben
  \item \textbf{: Bei den Berechtigungsanfragen während der Installation überlegen: "Braucht diese App das wirklich?"}Anhang B: Anregungen für praktische Übungen
  \item \textbf{Grundlagen-Übung}: Smartphone-Speicher aufräumen – Ziel ist es, vergessene alte Videos, doppelte Screenshots und lange ungenutzte Apps zu finden und zu löschen.
  \item \textbf{Software-Übung}: WeChat-Speicherbereinigung einrichten. Pfad: Ich Einstellungen Allgemein Speicherplatz.
\end{itemize}

\booksection{App-Installation – “Sicherheit zuerst”}

\begin{itemize}
  \item \textbf{: Akkugesundheit ansehen – iPhone unter Einstellungen Akku Akkugesundheit.}Allgemeine Fähigkeit
  \item \textbf{: Bei Problemen zuerst neu starten! Das ist der Universalschlüssel, der mehr als die Hälfte aller kleinen Smartphone-Wehwehchen löst.}Anhang C: Glossar der "Fachsprache"
\end{itemize}
